\section{Khảo sát bài toán xây dựng trợ lý ảo}
\subsection{Bối cảnh bài toán}

Trong môi trường nghiệp vụ của Ban Cơ yếu Chính phủ, nhu cầu tra cứu thông tin, hướng dẫn thủ tục, khai thác biểu mẫu và hỗ trợ hỏi đáp nội bộ ngày càng tăng. Người dùng thường cần tiếp cận nhanh các thông tin như quy trình xử lý hồ sơ, văn bản hướng dẫn, biểu mẫu nghiệp vụ, thông tin đào tạo, quy định nội bộ hoặc các nội dung hỗ trợ tác nghiệp thường gặp.

Hiện nay, các thông tin này thường được công bố hoặc lưu trữ trên nhiều kênh khác nhau như cổng thông tin, website đơn vị, tài liệu PDF, biểu mẫu điện tử, văn bản hướng dẫn, email hoặc các đầu mối hỗ trợ trực tiếp. Việc dữ liệu phân tán ở nhiều nguồn khiến người dùng mất thời gian tìm kiếm, đặc biệt khi không biết chính xác thông tin cần tra cứu nằm ở đâu hoặc được cập nhật trong văn bản nào.

Bên cạnh đó, trong quá trình vận hành thực tế, số lượng câu hỏi gửi đến bộ phận phụ trách có thể tăng mạnh và nhiều câu hỏi có nội dung lặp lại. Ví dụ như: ``Thủ tục xử lý hồ sơ này gồm những bước nào?'', ``Biểu mẫu cần dùng ở đâu?'', ``Văn bản hướng dẫn hiện hành là văn bản nào?'', ``Quy trình phê duyệt thực hiện ra sao?''. Nếu toàn bộ các câu hỏi này đều được xử lý thủ công, thời gian phản hồi sẽ chậm, chất lượng hỗ trợ khó đồng đều và gây áp lực cho cán bộ phụ trách.

Vì vậy, việc xây dựng một hệ thống trợ lý ảo ứng dụng trong Ban Cơ yếu Chính phủ có khả năng tiếp nhận câu hỏi tự nhiên, phân tích nội dung câu hỏi, truy xuất dữ liệu từ nguồn tri thức nghiệp vụ và đưa ra câu trả lời phù hợp là cần thiết. Hệ thống này giúp người dùng tiếp cận thông tin nhanh hơn, đồng thời hỗ trợ đơn vị giảm tải công việc hỏi đáp lặp lại và nâng cao hiệu quả phục vụ.

\begin{figure}[H]
    \centering
    \includegraphics[width=\textwidth]{content/source/Hình 2.1 Mô hình tư vấn tuyển sinh truyền thống và mô hình có chatbot hỗ trợ..png}
    \caption{Mô hình hỗ trợ thông tin truyền thống và mô hình có trợ lý ảo hỗ trợ}
    \label{fig:admission_consulting_model}
\end{figure}

\subsection{Thực trạng hỗ trợ thông tin hiện nay}

Qua khảo sát bài toán, hoạt động hỗ trợ người dùng hiện nay thường được triển khai theo các hình thức chính sau:

Thứ nhất, người dùng tự tìm kiếm thông tin trên cổng thông tin hoặc website chính thức của đơn vị. Đây là nguồn dữ liệu có độ tin cậy cao, tuy nhiên thông tin thường được trình bày dưới dạng bài viết dài, bảng biểu hoặc file PDF nên người dùng phải tự đọc và tự đối chiếu thông tin.

Thứ hai, người dùng đặt câu hỏi qua email, hotline hoặc đầu mối hỗ trợ trực tiếp. Hình thức này cho phép nhận được phản hồi từ cán bộ phụ trách, tuy nhiên khi số lượng yêu cầu lớn, thời gian phản hồi có thể bị kéo dài và các câu hỏi lặp lại làm tăng khối lượng công việc.

Thứ ba, một số nội dung được tổng hợp dưới dạng danh sách câu hỏi thường gặp hoặc tài liệu hướng dẫn. Cách này hỗ trợ tốt cho một số tình huống phổ biến, nhưng người dùng vẫn phải tự tìm đúng mục liên quan và khó sử dụng khi câu hỏi được diễn đạt theo nhiều cách khác nhau.

Thứ tư, một số chatbot truyền thống được xây dựng theo kịch bản cố định. Các chatbot này có thể xử lý các câu hỏi đơn giản nhưng thường gặp khó khăn với câu hỏi tự nhiên, câu hỏi thiếu ngữ cảnh hoặc câu hỏi yêu cầu trích xuất thông tin từ tài liệu nghiệp vụ dài.

Từ thực trạng trên, bài toán xây dựng trợ lý ảo cho Ban Cơ yếu Chính phủ cần một hệ thống có khả năng kết hợp giữa xử lý ngôn ngữ tự nhiên, quản lý hội thoại, truy xuất dữ liệu có cấu trúc và truy xuất tri thức từ tài liệu.

\begin{table}[H]
    \centering
    \caption{So sánh các hình thức hỗ trợ thông tin hiện nay}
    \label{tab:comparison_admission_consulting_methods}
    \begin{tabular}{|p{3.4cm}|p{5.2cm}|p{5.2cm}|}
        \hline
        \textbf{Hình thức hỗ trợ} & \textbf{Ưu điểm} & \textbf{Hạn chế} \\ \hline
        Website/cổng thông tin & Chính thống, đầy đủ thông tin & Khó tìm kiếm, nhiều tài liệu dài \\ \hline
        Email/đầu mối hỗ trợ & Có người hỗ trợ trực tiếp & Phản hồi chậm khi quá tải \\ \hline
        Hotline & Nhanh, trực tiếp & Phụ thuộc thời gian làm việc \\ \hline
        FAQ/tài liệu hướng dẫn & Dễ triển khai & Chỉ hỗ trợ câu hỏi tương đối cố định \\ \hline
        Chatbot kịch bản & Tự động hóa một phần & Khó xử lý câu hỏi tự nhiên \\ \hline
        Trợ lý ảo Rasa + RAG & Tự động, linh hoạt, truy xuất tài liệu & Cần xây dựng dữ liệu và quản trị tri thức \\ \hline
    \end{tabular}
\end{table}

\subsection{Đặc điểm của bài toán xây dựng trợ lý ảo}

Bài toán xây dựng trợ lý ảo ứng dụng trong Ban Cơ yếu Chính phủ có một số đặc điểm riêng biệt so với các bài toán hỏi đáp thông thường:

Thứ nhất, nguồn tri thức nghiệp vụ có thể thay đổi theo thời gian và phụ thuộc vào văn bản, quy định hoặc quy trình đang có hiệu lực. Vì vậy, hệ thống cần hỗ trợ cập nhật dữ liệu linh hoạt và có cơ chế quản lý phiên bản hoặc trạng thái hiệu lực của tài liệu.

Thứ hai, câu hỏi của người dùng thường ngắn, thiếu ngữ cảnh và được diễn đạt bằng ngôn ngữ tự nhiên. Ví dụ: ``Biểu mẫu này lấy ở đâu?'', ``Quy trình xử lý hồ sơ gồm mấy bước?'', ``Văn bản hướng dẫn hiện hành là văn bản nào?''. Hệ thống cần xác định được ý định của người dùng và các thông tin còn thiếu để hỏi lại khi cần.

Thứ ba, người dùng có thể sử dụng từ viết tắt, cách gọi không chính thức hoặc cách diễn đạt khác nhau cho cùng một nội dung. Do đó, hệ thống cần có khả năng chuẩn hóa các cách diễn đạt này về dữ liệu chính thức trong hệ thống.

Thứ tư, nhiều câu hỏi cần truy xuất dữ liệu chính xác từ tài liệu hoặc cơ sở dữ liệu. Các câu hỏi như ``Biểu mẫu áp dụng cho thủ tục này là gì?'' hoặc ``Quy trình phê duyệt hiện hành được quy định ở văn bản nào?'' cần được trả lời dựa trên nguồn dữ liệu chính thống, không nên trả lời theo suy đoán.

Thứ năm, một số câu hỏi mang tính hướng dẫn hoặc gợi ý thao tác. Ví dụ: ``Tôi cần bắt đầu từ bước nào?'', ``Nếu thiếu giấy tờ thì xử lý ra sao?''. Các câu hỏi này đòi hỏi hệ thống không chỉ truy xuất dữ liệu mà còn phải tổ chức lại thông tin theo hướng dễ hiểu, phù hợp với ngữ cảnh sử dụng.

\begin{table}[H]
    \centering
    \caption{Một số ví dụ câu hỏi nghiệp vụ và dạng xử lý tương ứng}
    \label{tab:admission_question_examples}
    \begin{tabular}{|p{4.1cm}|p{3.2cm}|p{3.1cm}|p{3.8cm}|}
        \hline
        \textbf{Câu hỏi người dùng} & \textbf{Ý định} & \textbf{Thực thể cần nhận diện} & \textbf{Hướng xử lý} \\ \hline
        Biểu mẫu cho thủ tục này ở đâu? & Hỏi biểu mẫu & Tên thủ tục & Tra cứu tài liệu, biểu mẫu \\ \hline
        Quy trình xử lý hồ sơ gồm những bước nào? & Hỏi quy trình & Loại hồ sơ & Truy xuất quy trình nghiệp vụ \\ \hline
        Văn bản hướng dẫn hiện hành là gì? & Hỏi văn bản & Nhóm nghiệp vụ & Tra cứu văn bản hiệu lực \\ \hline
        Nếu thiếu giấy tờ thì xử lý thế nào? & Hỏi hướng dẫn xử lý & Loại giấy tờ, thủ tục & Hướng dẫn theo quy định \\ \hline
        Tôi cần liên hệ bộ phận nào? & Hỏi đầu mối hỗ trợ & Nhóm nghiệp vụ & Tra cứu thông tin liên hệ \\ \hline
    \end{tabular}
\end{table}

\subsection{Mục tiêu xây dựng hệ thống}

Mục tiêu của hệ thống trợ lý ảo ứng dụng trong Ban Cơ yếu Chính phủ là xây dựng một công cụ hỗ trợ tự động, giúp người dùng tra cứu và khai thác thông tin nghiệp vụ một cách nhanh chóng, chính xác và thuận tiện.

Các mục tiêu chính gồm:
\begin{itemize}
    \item Hỗ trợ người dùng đặt câu hỏi bằng ngôn ngữ tự nhiên, không cần thao tác qua nhiều menu phức tạp.
    \item Tự động nhận diện ý định và thông tin quan trọng trong câu hỏi.
    \item Cung cấp câu trả lời dựa trên dữ liệu và tài liệu nghiệp vụ chính thống.
    \item Tích hợp hệ thống RAG để truy xuất thông tin từ tài liệu dài như văn bản hướng dẫn, quy định, biểu mẫu hoặc file PDF nội bộ.
    \item Cung cấp phân hệ quản trị để cập nhật và kiểm soát dữ liệu, tài liệu RAG, lịch sử hỏi đáp và cấu hình phục vụ chatbot.
    \item Nâng cao chất lượng hỗ trợ và giảm tải cho cán bộ phụ trách thông qua tự động hóa các câu hỏi phổ biến.
\end{itemize}

\begin{figure}[H]
    \centering
    \includegraphics[width=\textwidth]{content/source/Hình 2.2 Mục tiêu tổng quát của hệ thống chatbot tuyển sinh..png}
    \caption{Mục tiêu tổng quát của hệ thống trợ lý ảo}
    \label{fig:chatbot_system_objectives}
\end{figure}
